\section{Poison}\label{Poison}
  Poisons are organic substances that are dangerous to living creatures.
  They can deal damage or inflict debilitating effects.
  Some effects which are not literally poisonous, such as animal venom or fungal spores, are considered poisons.
  Poisons are not \glossterm{conditions}, and cannot be removed by abilities that remove conditions (see \pcref{Conditions}).
  Common poisons are listed in \tref{Consumables}.
  You can use the Craft (poison) skill to create poisons (see \pcref{Crafting Items}).

  Poisons always have the \atPoison tag, even if the ability inflicting the poison does not have that tag.
  They never apply \glossterm{extra damage} from any source.

  \subsection{Poison Effects}\label{Poison Effects}
    When you come in contact with a poison, it makes a \glossterm{reactive attack} against your Fortitude defense.
    On a hit, you become \glossterm{poisoned} by the poison.
    On a miss, the poison has no effect on you.

    While a creature is poisoned, the poison makes an attack against its Fortitude at the end of its turn.
    On a hit, the poison escalates.
    Some poisons have stronger effects after they escalate a certain number of times.
    On a miss, the creature makes progress towards removing the poison.
    If the poison misses three times, the creature stops being poisoned.

    If you become poisoned again by the same poison, it does not intensify the effects of the poison.
    However, it cancels any progress you had made towards removing the poison.
    A poison is considered the same if it has the same name.
    If you are affected by multiple different poisons with the same name, but different accuracy bonuses, use the highest accuracy bonus.

    Poisons have no effect on non-living creatures.

    \subsubsection{Poison Stages}
      Poison effects are divided into stages.
      Becoming poisoned does not have any ill effects until the poison progresses to its first stage.

      Many poisons have a Stage 1 effect.
      This effect happens as soon as the poison's attack first succeeds against you.
      Some poisons also have a Stage 3 effect.

    \subsubsection{Poison Accuracy}
      A poison's accuracy depends on the way it was applied.
      Item-based poisons have a specific accuracy listed in their description.
      This accuracy does not depend on the skill of the creature inflicting the poison.

      Poisons inflicted by creature abilities use the creature's \glossterm{accuracy}.
      They may also have additional modifiers listed in the ability's description.
      For monsters, the poison's accuracy will be listed in the monster description.

  \subsection{Poison Transmission}\label{Poison Transmission}\label{Transmission}

    There are three ways that poisons can be contracted.

    \parhead{Contact} A contact poison affects any creature that touches it with bare skin.
    \parhead{Ingestion} An ingestion poison affects any creature that eats, drinks, or breathes it, depending on the type of poison.
    Ingestion poisons have no effect when touched or used to coat weapons.
    This means you can avoid negative effects from ingestion poisons by holding your breath and avoiding eating or drinking.
    \parhead{Injury} An injury poison affects any creature that suffers an \glossterm{injury} from something bearing the poison.
    Almost all injury poisons take liquid form, and are typically used to coat weapons.

  \subsection{Poison Forms}\label{Poison Forms}

    There are four forms of poison.
    If a poison can be thrown, throwing one dose of the poison requires a standard action.

    \parhead{Gas} Gaseous poisons are difficult to store, but easy to affect foes with.
    Unless otherwise noted, a gas poison can be thrown at a location within \shortrange, and affects a \tinyarea radius \glossterm{zone}.
    Once active, gas poisons typically linger for one minute before dispersing, though high wind can clear them much more rapidly.
    Walking through the gas cloud will affect creatures with contact poisons, but generally not with ingestion poisons.
    \parhead{Liquid} Liquid poisons are the most common type of poison.
    Liquid poisons can be used to coat weapons, slipped into food, or simply thrown at foes.
    A dose of a liquid poison is usually about one ounce of the poison.
    Unless otherwise noted, a liquid poison can be thrown at something within \shortrange.
    \parhead{Pellet} Some rare poisons come in small, solid pellets or cubes.
    Typically, these pellets contain a powerful liquid poison that becomes inert quickly after being exposed.
    Pellet poisons are typically applied by being slipped into food.
    \parhead{Powder} Poison in powder form cannot be used to coat weapons, but can be slipped into food or thrown at foes.
    % adjacent or 10 feet?
    Unless otherwise noted, a powder poison can be thrown at something within 10 foot range.

  \subsection{Coating Weapons with Poison}\label{Coating Weapons with Poison}
    As a standard action, you can coat a weapon with a single dose of a liquid contact-based or injury-based poison.
    The next time a creature takes damage from a \glossterm{strike} using that weapon, the struck creature comes in contact with the poison.
    This removes one dose of the poison from the weapon.
    Coated poisons expire and lose their effectiveness after ten minutes.

    An injury-based poison has no effect if the strike does not cause an \glossterm{injury}, but the dose is still removed from the weapon.
    For this reason, injury-based poisons are typically applied to secondary weapons that can be used after the subject is already weakened.

    A weapon can hold up to three poison doses of the same poison.
    Mixing different poison types on the same weapon is ineffective, as each poison dilutes the others.
    Only the highest rank poison on the weapon has any effect.

  \subsection{Poison Materials}\label{Poison Materials}
    Creating a poison requires special materials.
    The type of materials required, and how those materials can be acquired, depend on the type of poison:

    \begin{raggeditemize}
      \itemhead{Alchemical} Alchemical poisons require alchemical materials.
        These normally can't be found in nature.
        In unusual circumstances, these components can be synthesized from natural chemicals or magical materials.
        This takes an hour of work, and requires a Craft (alchemy) check equal to 10 \add twice the poison's \glossterm{rank}.
      \itemhead{Plant} Plant-based poisons can typically be harvested by making a Survival \glossterm{extended check} to search in appropriate terrain.
        This usually takes an hour of work.
        The \glossterm{difficulty value} is usually equal to 10 \add twice the poison's \glossterm{rank}.
      \itemhead{Venom} Venom requires an appropriate body part from a creature -- often, poison it naturally produces.
        This usually takes a minute of work.
        Harvesting a body part from a creature requires a Medicine or Survival check equal to 5 \add twice the poison's rank.
    \end{raggeditemize}
